\uuid{cdjy}
\exo7id{259}
\titre{exo7 259}
\auteur{bodin}
\organisation{exo7}
\datecreate{1998-09-01}
\isIndication{false}
\isCorrection{false}
\chapitre{Arithmétique dans Z}
\sousChapitre{Divisibilité, division euclidienne}
\cpge{mpsi-mp2i.arithmetique-entiers.d, mpsi-mp2i.calcul-algebrique-trigonometrie.a, pcsi.calcul-algebrique-trigonometrie.a, ptsi.calcul-algebrique-trigonometrie.a, pcsi.raisonnement-vocabulaire-ensembliste.c, ptsi.raisonnement-vocabulaire-ensembliste.c}
\module{Algèbre}
\niveau{L1}
\difficulte{}

\contenu{
\question{
On consid\`ere le nombre $m = 2^np$, dans lequel $n$ d\'esigne
un entier naturel quelconque et $p$ un nombre premier. Dresser
la liste des diviseurs de $m$, y compris $1$ et $m$ lui-m\^eme,
et calculer, en fonction de $m$ et $p$, la somme $S$ de tous
ces diviseurs.
}
}
